\section{Discussion}
\label{sec:discussion}
The question this article set out to answer was never simply whether a prompt is useful in general. The real question is: once a coarse box has already narrowed the search, does the way the model handles that prompt still matter? Attention U-Net without a box answers a different question altogether, and its low Dice, near 0.34 on FracAtlas, measures only how much of the task is finding the lesion. The design is tested wherever the same box is given to another model or stripped of a component: the prompt-matched Attention U-Net baselines, the fine-tuned promptable foundation models, and the ablation. Across all of them, PGA-UNet keeps a margin that is modest on BTXRD, clearer on FracAtlas, and largest on the Attention U-Net bottom-Dice group and on the smallest lesions. If the point of keeping the prompt active through the encoder and decoder is to preserve a faint lesion signal, it should help most where that signal is faintest, and that is the direction the results take: the margin over prompt-matched baselines, the per-component ablation gaps, and the robustness to prompt displacement are all largest on FracAtlas, on the smallest lesions, and under off-center prompts, and all shrink toward the easy end. The consistency of that direction across independent measurements is the main support for the design, more than any single comparison.

Beyond the main result, the study also yields two side findings. First, the loss comparison gives a fairly clear negative result: neither the size-conditioned Tversky loss nor the Focal Dice loss helps, and Focal Dice even makes things distinctly worse on FracAtlas, so we keep the default Dice + BCE objective. Second, the self-assessment score gives a moderate positive result: $Q$ correlates with the true Dice at a rank level of about 0.5 to 0.7 without needing any label at all, and it succeeds where a small learned regression head does not, since that head ultimately just predicts a mean value. Even so, $Q$ remains only a heuristic number, not yet calibrated, so the candidate list built from it should only be used to suggest where a clinician might look first, not to replace self-detection.

This study also has fairly concrete limitations. Inputs are resized to a fixed square, and the resolution results themselves show that this genuinely costs accuracy on the smallest lesions. The prompts in this article are all drawn from existing annotations, not directly drawn by a radiologist, and the protocol assumes the user has already found the correct region before drawing a box over it. Cases where the box only partially covers the lesion, is negative, or targets the wrong region are not addressed by this article; a direct check shows that the self-assessment score $Q$ in particular stays high on many boxes drawn over lesion-free anatomy, because the model's weights were learned only on lesion-bearing images. Making $Q$ also signal lesion absence would require training with lesion-free regions, or pairing $Q$ with a separate lesion-presence gate. BTXRD and FracAtlas are trained completely separately, with no joint model and no cross-dataset test step, and patient-level separation cannot be confirmed from the public metadata. The Monte Carlo runs only show stability across splits, not superiority. The per-component ablation is on a single split, so it shows the direction of each component's effect but cannot settle the finer contributions; a multi-split run is future work. Finally, the gains under prompt perturbation are consistent with the design intent but have not been verified by direct feature-level analysis, so the mechanism is inferred from behavior rather than measured inside the network.
